%%%%%%%%%%%%%%%%%%%%%%%%%%%%%%%%%%%%%%%%
\chapter{Introducción}
%%%%%%%%%%%%%%%%%%%%%%%%%%%%%%%%%%%%%%%%

En los últimos años, los grandes modelos de lenguaje (\textit{Large Language Models}, LLMs) han transformado múltiples áreas mediante avances en procesamiento de lenguaje natural (\textit{Natural Language Processing}, NLP). Modelos como GPT-4, LLaMA y Mistral han demostrado capacidades emergentes en tareas de comprensión y generación de lenguaje, permitiendo aplicaciones que van desde asistentes conversacionales hasta análisis de documentos especializados y generación de código.

El rendimiento de estos modelos no depende únicamente de su arquitectura o datos de pre-entrenamiento, sino también de las técnicas de \textit{fine-tuning} (ajuste fino) empleadas para adaptarlos a dominios específicos. El paradigma de pre-entrenar y ajustar modelos de lenguaje a tareas específicas se ha establecido como el enfoque estándar en el área \parencite[p.~220]{Ding2023}, permitiendo que modelos generales se especialicen mediante entrenamiento adicional en datos de dominio. Este principio de especialización se fundamenta en evidencia empírica que demuestra que los modelos adaptados mediante \textit{fine-tuning} pueden capturar mejor el conocimiento específico de la tarea, conduciendo a rendimiento superior en dominios restringidos comparado con modelos de propósito general.

Desde 2020 han surgido múltiples técnicas representativas de \textit{fine-tuning}, cada una con enfoques distintos de adaptación. En el marco de esta revisión, se analizan nueve técnicas representativas organizadas en tres familias según su paradigma: \textit{Full Fine-Tuning} y su variante \textit{Supervised Fine-Tuning} (SFT), que representan el enfoque tradicional de actualización completa de parámetros sobre datos etiquetados; los métodos de \textit{Parameter-Efficient Fine-Tuning} (PEFT), como LoRA, QLoRA, \textit{Adapter Layers}, \textit{Prefix Tuning} y \textit{Prompt Tuning}, que priorizan eficiencia reduciendo drásticamente el número de parámetros entrenables; y los métodos de alineación (\textit{Alignment Methods}), como RLHF y DPO, que ajustan el comportamiento del modelo para reflejar preferencias y valores humanos. La coexistencia de estos tres paradigmas en la literatura reciente plantea preguntas sobre cuándo cada enfoque resulta más adecuado según el contexto operativo, pregunta central que esta revisión busca responder con base en evidencia empírica.

La evidencia empírica sobre estas técnicas se encuentra dispersa en múltiples fuentes académicas, donde cada publicación reporta comparaciones o evaluaciones según objetivos específicos del estudio. La proliferación de métodos ha sido notable: en los últimos años se han publicado más de cien artículos sobre PEFT \parencite[p.~2]{Lialin2024}, reflejando el intenso interés de la comunidad investigadora. Algunos trabajos comparan técnicas dentro de un paradigma, otros evalúan una técnica en múltiples dominios de aplicación, mientras que otros se enfocan en dominios particulares como procesamiento de textos médicos, generación de código o análisis de documentos legales. Esta especialización natural, aunque valiosa para avanzar conocimiento en áreas específicas, genera evidencia fragmentada que dificulta la consolidación de conocimiento comparativo entre técnicas y entre dominios.

La comparabilidad directa de resultados entre estudios se ve limitada por heterogeneidad en diseños experimentales. Estudios comparativos recientes evidencian que incluso los artículos más representativos carecen de comparabilidad completa, debido a diferencias en configuraciones de evaluación como conteos de parámetros variables o ausencia de ciertas métricas \parencite[p.~35]{Lialin2024}. La variabilidad en resultados reportados es sustancial: estudios experimentales documentan amplios rangos de rendimiento según la combinación de técnica, escala de modelo y configuración experimental empleada \parencite[p.~29, Tabla~4]{Lialin2024}. Adicionalmente, la combinación óptima de técnicas PEFT ``puede variar para diferentes \textit{backbones} de PLM, tareas \textit{downstream} y escalas de datos'' \parencite[p.~221]{Ding2023}, demostrando que el rendimiento depende críticamente del contexto experimental específico. En ausencia de síntesis comparativa sistemática, la selección de técnicas frecuentemente se basa en disponibilidad de implementaciones públicas o recomendaciones anecdóticas, más que en evidencia empírica consolidada.

Una revisión sistemática empleando el protocolo PRISMA (\textit{Preferred Reporting Items for Systematic Reviews and Meta-Analyses}) aborda estas limitaciones mediante un proceso estructurado y transparente que documenta explícitamente la estrategia de búsqueda, establece criterios claros de inclusión y exclusión, y registra sistemáticamente el proceso de selección y extracción de datos. Esta sistematicidad permite consolidar evidencia dispersa, identificar patrones recurrentes en rendimiento según dominios de aplicación y detectar vacíos en comparaciones existentes. La propia literatura reconoce esta necesidad: la falta de comparabilidad entre estudios y la ausencia de protocolos de evaluación unificados constituyen limitaciones recurrentes en el campo \parencite[p.~35]{Lialin2024}. Desde la perspectiva de la Ciencia de Datos, esta revisión se enfoca en métricas de calidad de predicción y patrones de aplicabilidad según dominio, excluyendo deliberadamente el análisis de eficiencia computacional dado que la heterogeneidad de implementaciones y hardware impediría comparaciones válidas entre estudios.

El presente trabajo aborda la siguiente pregunta de investigación:
\begin{quote}
¿Cuáles son las características, el rendimiento reportado en distintos dominios de aplicación, y los patrones de aplicabilidad de las técnicas de \textit{fine-tuning} para \textit{Large Language Models} desarrolladas entre 2020 y 2025, según la evidencia reportada en la literatura académica?
\end{quote}

El alcance de esta revisión se delimita de la siguiente manera. Se analizan nueve técnicas representativas que cubren los tres paradigmas identificados: \textit{Full Fine-Tuning} y SFT como enfoques de actualización completa de parámetros, cinco técnicas de \textit{Parameter-Efficient Fine-Tuning} (LoRA, QLoRA, \textit{Adapter Layers}, \textit{Prefix Tuning}, \textit{Prompt Tuning}) y dos métodos de alineación (RLHF, DPO). La revisión cubre literatura publicada entre 2020 y 2025, período que corresponde a la explosión metodológica en técnicas de \textit{fine-tuning} tras la introducción de modelos de lenguaje de gran escala. Se incluyen artículos de revistas académicas con evaluación empírica cuantitativa, publicados en inglés y con acceso abierto, asegurando reproducibilidad del proceso de selección. El análisis se organiza en torno a cuatro dominios de aplicación: Salud y Biomedicina, Lenguaje y Humanidades, Ciencias Exactas y Computación, e Industria y Negocios. La aplicación de estos criterios mediante metodología PRISMA resultó en la inclusión de 84 estudios empíricos, cuya evidencia constituye la base del análisis comparativo presentado en esta monografía.